% glyphs.tex — the emoji substitution table.
%
% Moved out of book-header.tex unchanged at improvement #79; extended at #81.
%
% Tectonic's fonts have no emoji glyphs, so the book's callout vocabulary (BOOK-SPEC.md
% § 7 and the write-topic-docs skill) prints as tofu boxes without this. Each of the six
% sanctioned marks is mapped to a typographic equivalent that carries the same meaning in
% print. EPUB and HTML keep the real emoji — they render fine there.
%
% Since #81, scripts/lua/callouts.lua consumes 💡, 🔑 and the ⚠️ that opens a gotcha
% before typesetting starts: the mark becomes a block with an uppercase label, which says
% the same thing louder. What is left for this table is every *other* appearance — the
% vocabulary table in the front matter, a ⚠️ inside a list item, a 💡 quoted in prose —
% and the two dingbats, which the filter never touches.

\usepackage{newunicodechar}
\usepackage{pifont}

% ✅ / ❌ — these two carry meaning inside tables and before/after pairs, so they
% become real dingbats rather than words. They work in one ink because the two glyphs
% differ in shape rather than in colour.
\newunicodechar{✅}{\ding{51}}
\newunicodechar{❌}{\ding{55}}

% ⚠️ — a gotcha. A boxed bang reads as a warning at body size and, unlike a bold `!` on
% its own, does not disappear into the sentence around it.
\newcommand{\bookwarnglyph}{%
  {\setlength{\fboxsep}{1.2pt}\setlength{\fboxrule}{\tokHairline}%
   \fbox{\labelfont\fontsize{\tokMicroSize}{\tokMicroSize}\selectfont !}}%
}
\newunicodechar{⚠}{\bookwarnglyph}

% 💡 and 🔑 — a solid square, sized in em so it takes the weight and the size of whatever
% it is set in rather than arriving as a dingbat from a third face. #79 dropped both
% silently, which left the front matter's own vocabulary table describing two marks it
% did not print.
\newcommand{\booksquareglyph}{\raisebox{0.06em}{\rule{0.5em}{0.5em}}}
\newunicodechar{💡}{\booksquareglyph}
\newunicodechar{🔑}{\booksquareglyph}

% ---------------------------------------------------------------------------
% The four CJK characters — improvement #93
% ---------------------------------------------------------------------------
%
% 年 月 日 ￥, twelve occurrences, all of them inside the `Intl` samples in
% Frontend/Internationalization/03-date-number-formatting.md — a comment showing what
% `Intl.DateTimeFormat('ja-JP')` actually returns. #77 measured them printing as holes
% and recorded the fix rather than taking it, because it needed a font download and a
% licence call.
%
% Both are now made. `assets/fonts/NotoSansJP-Subset.otf` is **3.9 KB**: Noto Sans JP,
% OFL 1.1 — the same licence as the three Source families — subset with `fontTools` to
% exactly these four glyphs. #81's objection, that a CJK face costs "an order of
% magnitude more than all three Source families", is true of a whole face and not of four
% characters.
%
% 🔴 The replacement text must use \char rather than the character itself.
% \newunicodechar makes the character active, so writing 年 inside its own definition
% recurses until TeX runs out of stack.

\newfontfamily\bookcjkfont{NotoSansJP-Subset}[
  Path = \bookfontpath,
  Extension = .otf,
  Scale = MatchLowercase,
]

\newunicodechar{年}{{\bookcjkfont\char"5E74}}
\newunicodechar{月}{{\bookcjkfont\char"6708}}
\newunicodechar{日}{{\bookcjkfont\char"65E5}}
\newunicodechar{￥}{{\bookcjkfont\char"FFE5}}

% ★ — U+2605, inside the rating-stars sample in the web-components chapter. It is what the
% code renders, so it has to print rather than be reworded. Zapf Dingbats has it. #93.
\newunicodechar{★}{\ding{72}}

% U+FE0F, the variation selector that trails ⚠️ and friends in the source.
\newunicodechar{️}{}

% ---------------------------------------------------------------------------
% Box-drawing characters
% ---------------------------------------------------------------------------
%
% #81 was written expecting to map these: "117 files still carry box-drawing characters
% that print as tofu". Measured against the manuscript, that is wrong twice over.
%
% There are five of them — ─ │ ├ └ █ — across **8 files, 175 occurrences, and every
% single one is inside a code fence.** Source Code Pro carries all five, so they already
% print correctly; only Source Serif 4 and Source Sans 3 are missing them, and no
% box-drawing character appears in prose or in a table anywhere in the book.
%
% Substituting them here would therefore make the book worse, not better: a \rule wide
% enough to read is not the monospaced width the ASCII diagram was drawn to, so every
% tree and every timeline in those 8 files would lose its alignment to fix a problem it
% does not have. They are deliberately left alone. The missing-character count that
% build-book.sh reports is what catches it if one ever moves into prose.
%
% One character genuinely is missing from all three families: ✕ U+2715, which appears
% once, as the label of a close button in an HTML sample.
\newunicodechar{✕}{\ding{56}}

% ---------------------------------------------------------------------------
% What the build's missing-character report actually found
% ---------------------------------------------------------------------------
%
% #81 added the report to build-book.sh and then ran it over the whole manuscript, which
% is the only way any of this was ever going to be known. It found six characters the
% vendored faces cannot set — none of them a box-drawing character.
%
% α U+03B1 is the inverse Ackermann function in the Union-Find complexity, and it is a
% mathematical symbol rather than a Greek letter in running text, so it is set as one.
% latinmodern-math is already loaded by pandoc's own template.
\newunicodechar{α}{\ensuremath{\alpha}}

% The other four are 年 月 日 and ￥, twelve occurrences between them, all inside the
% `Intl.DateTimeFormat` and `Intl.NumberFormat` samples in the internationalisation
% chapter. They are **deliberately not substituted**, because there is no honest
% substitute: those comments show what the ja-JP locale actually returns, and a
% romanisation would make the code sample claim an output it does not produce.
%
% Setting them needs a CJK face vendored beside the other three, which is a licence and
% a file-size decision of the same shape as #79's and costs an order of magnitude more
% than all three Source families together. Recorded against #77, which owns the final
% proof. Until then the build says so on every run rather than printing a hole quietly.
